\documentclass[fontsize=10pt,oneside,parskip=half]{scrreprt}
\usepackage{ulille-rapport}
\usepackage{amsmath}
\usepackage{float}
\usepackage{booktabs}
\usepackage{tabularx}
\usepackage{listings}

\newcommand{\ic}[1]{{\setlength{\fboxsep}{2pt}\colorbox{BeigeSerein}{\texttt{#1}}}}

\subject{SAE S2.02 : Rapport graphes}
\title{Première Version}
\author{Yann Renard \\ Tom Cox \\ Alexandre}
\direction{Équipe E5}

\lstdefinestyle{java}{
  language=Java,
  basicstyle=\ttfamily\small,
  keywordstyle=\bfseries\color{IndigoInclusif},
  commentstyle=\itshape\color{GrisStatutaire},
  stringstyle=\color{VertAvenir},
  backgroundcolor=\color{BeigeSerein},
  frame=single,
  framesep=6pt,
  rulecolor=\color{BleuHorizon},
  breaklines=true,
  breakatwhitespace=true,
  showstringspaces=false,
  tabsize=2,
  numbers=left,
  numberstyle=\tiny\color{GrisStatutaire},
  numbersep=8pt,
  xleftmargin=20pt,
  xrightmargin=4pt,
}

\begin{document}
\selectlanguage{french}

\maketitle
\pagecolor{Blanc}
{%
  \setlength{\parskip}{0pt}%
  \tableofcontents%
}

% ─────────────────────────────────────────────
\chapter{Version 1 : Un seul moyen de transport}
% ─────────────────────────────────────────────

% ──────────────────────────────
\section{Exemple}
% ──────────────────────────────

Pour illustrer la Version~1, nous avons choisi un réseau de transport reliant
cinq villes du nord de la France : Lille, Dunkerque, Cambrai, Paris et
Crépy-en-Valois. Ce réseau comprend dix-huit connexions réparties sur trois modes
de transport (train, bus et avion) avec des données de coûts estimées à partir
des vitesses moyennes de chaque mode et des barèmes d'émissions carbone de
l'ADEME.

\noindent\textbf{Tableau 1 : Données complètes du réseau de transport}

\begin{table}[H]
  \centering
  \begin{tabular}{lllrrr}
    \toprule
    \textbf{Départ} & \textbf{Arrivée} & \textbf{Mode}
      & \textbf{Prix (€)} & \textbf{CO\textsubscript{2} (kg)}
      & \textbf{Durée (min)} \\
    \midrule
    Lille             & Dunkerque         & TRAIN  & 18  & 0{,}8  & 65  \\
    Dunkerque         & Lille             & TRAIN  & 18  & 0{,}8  & 65  \\
    Lille             & Cambrai           & TRAIN  & 14  & 0{,}7  & 55  \\
    Cambrai           & Lille             & TRAIN  & 14  & 0{,}7  & 55  \\
    Cambrai           & Paris             & TRAIN  & 35  & 1{,}8  & 95  \\
    Paris             & Cambrai           & TRAIN  & 35  & 1{,}8  & 95  \\
    Paris             & Crépy-en-Valois   & TRAIN  & 12  & 0{,}5  & 45  \\
    Crépy-en-Valois   & Paris             & TRAIN  & 12  & 0{,}5  & 45  \\
    Lille             & Paris             & TRAIN  & 55  & 2{,}4  & 60  \\
    Paris             & Lille             & TRAIN  & 55  & 2{,}4  & 60  \\
    Dunkerque         & Cambrai           & TRAIN  & 22  & 1{,}0  & 80  \\
    Cambrai           & Dunkerque         & TRAIN  & 22  & 1{,}0  & 80  \\
    Dunkerque         & Crépy-en-Valois   & TRAIN  & 45  & 2{,}0  & 150 \\
    Dunkerque         & Cambrai           & BUS    & 12  & 2{,}1  & 110 \\
    Cambrai           & Crépy-en-Valois   & BUS    & 18  & 2{,}8  & 130 \\
    Paris             & Crépy-en-Valois   & BUS    & 8   & 1{,}2  & 75  \\
    Lille             & Paris             & AVION  & 120 & 85{,}0 & 55  \\
    Dunkerque         & Paris             & AVION  & 110 & 80{,}0 & 50  \\
    \bottomrule
  \end{tabular}
  \caption{Données complètes du réseau de transport}
\end{table}

{\small\itshape Les durées intègrent les temps d'attente habituels
(enregistrement aéroport pour l'avion, correspondances en gare).}

\newpage
Le scénario retenu pour ce prototype est le suivant :

\begin{itemize}
  \item Mode de transport : TRAIN uniquement
  \item Critère d'optimisation : le prix (on cherche à minimiser le coût en euros)
  \item Ville de départ : Lille
  \item Ville d'arrivée : Crépy-en-Valois
  \item Nombre d'itinéraires souhaités : 4
\end{itemize}

Ce scénario est représentatif d'un voyageur contraint par son budget, souhaitant
rejoindre Crépy-en-Valois depuis Lille en train. Le Tableau~2 présente les quatre
meilleurs itinéraires identifiés, classés par ordre croissant de prix.

\noindent\textbf{Tableau 2 : Les 4 meilleurs itinéraires TRAIN de Lille à
Crépy-en-Valois (critère : prix)}

\begin{table}[H]
  \centering
  \begin{tabularx}{\textwidth}{Xrr}
    \toprule
    \textbf{Itinéraire} & \textbf{Prix (€)} & \textbf{Classement} \\
    \midrule
    Lille → Cambrai → Paris → Crépy-en-Valois      & 61 & 1\textsuperscript{er} \\
    Lille → Dunkerque → Crépy-en-Valois            & 63 & 2\textsuperscript{e}  \\
    Lille → Paris → Crépy-en-Valois                & 67 & 3\textsuperscript{e}  \\
    Lille → Cambrai → Dunkerque → Crépy-en-Valois  & 81 & 4\textsuperscript{e}  \\
    \bottomrule
  \end{tabularx}
  \caption{Les 4 meilleurs itinéraires TRAIN de Lille à Crépy-en-Valois (critère : prix)}
\end{table}

{\small\itshape Les itinéraires empruntant l'avion ou le bus sont exclus du
scénario, le filtre portant exclusivement sur le mode TRAIN.}

% ──────────────────────────────
\section{Modèle pour l'exemple}
% ──────────────────────────────

Pour résoudre ce problème, on modélise le réseau filtré par un graphe orienté
valué. Chaque ville du réseau devient un sommet, et chaque connexion TRAIN devient
une arête orientée dont le poids correspond au critère choisi, ici le prix en euros.

\begin{figure}[H]
  \centering
  \includegraphics[width=0.60\textwidth]{imgrapport/graphesv1}
  \caption{Graphe orienté valué, scénario TRAIN, critère PRIX (Lille vers Crépy-en-Valois)}
\end{figure}

Dans ce graphe, les cinq sommets sont les villes
\{Lille, Dunkerque, Cambrai, Paris, Crépy-en-Valois\}. Les arêtes orientées et
valuées représentent les connexions TRAIN disponibles. À titre d'exemple,
l'arête Lille~→~Paris porte le poids~55 (prix en euros), l'arête
Lille~→~Cambrai porte le poids~14, et l'arête Paris~→~Crépy-en-Valois porte le
poids~12.

Les quatre meilleurs chemins dans ce graphe, exprimés sous forme de suites
d'arêtes avec leur poids cumulé, sont :

\begin{itemize}
  \item \textbf{Chemin~1 :} Lille → Cambrai → Paris → Crépy-en-Valois,
        poids total $= 14 + 35 + 12 = \mathbf{61}$~€
  \item \textbf{Chemin~2 :} Lille → Dunkerque → Crépy-en-Valois,
        poids total $= 18 + 45 = \mathbf{63}$~€
  \item \textbf{Chemin~3 :} Lille → Paris → Crépy-en-Valois,
        poids total $= 55 + 12 = \mathbf{67}$~€
  \item \textbf{Chemin~4 :} Lille → Cambrai → Dunkerque → Crépy-en-Valois,
        poids total $= 14 + 22 + 45 = \mathbf{81}$~€
\end{itemize}

L'algorithme KPCC explore les chemins sans cycle dans le graphe et les retourne
ordonnés par poids croissant. Les résultats ci-dessus ont été obtenus et vérifiés
par l'exécution de la classe \ic{GraphesV1} décrite en Section~1.4.

% ──────────────────────────────
\section{Modélisation pour la Version~1 dans le cas général}
% ──────────────────────────────

On se donne un problème de recherche d'itinéraire défini par : un ensemble de
villes~$V$, un ensemble de connexions~$C$ entre ces villes, une modalité de
transport~$m$ choisie, un critère d'optimisation~$c$ parmi
\{TEMPS, PRIX, CO2\}, une ville de départ~$s$, une ville d'arrivée~$t$, et un
entier~$k$ représentant le nombre d'itinéraires souhaités. Voici comment ce
problème se traduit en termes de graphes.

\noindent\textbf{Type de graphe}

On utilise un \emph{multigraphe orienté valué}. L'orientation est nécessaire car
une connexion de la ville~$x$ vers la ville~$y$ n'implique pas l'existence d'une
connexion en sens inverse : les lignes de transport sont directionnelles. La
valuation permet d'associer à chaque arête le coût correspondant au critère
choisi. Enfin, la structure de multigraphe est justifiée par le fait que deux
villes pourraient théoriquement être reliées par plusieurs connexions de même
modalité, ce qui correspondrait à des arêtes parallèles dans le graphe.

\noindent\textbf{Sommets du graphe}

Les sommets du graphe sont les villes qui apparaissent dans au moins une
connexion de la modalité choisie, soit en départ, soit en arrivée. Autrement
dit, on ne retient que les villes réellement desservies par le mode de transport
sélectionné ; les autres sont ignorées. Dans notre exemple, filtrer sur TRAIN
exclut toute ville uniquement accessible en bus ou en avion.

\noindent\textbf{Arêtes du graphe et leurs poids}

Il existe une arête orientée du sommet~$x$ vers le sommet~$y$ dans le graphe
si et seulement si il existe une connexion de modalité~$m$ allant de~$x$ à~$y$
dans l'ensemble~$C$. Le poids de cette arête est la valeur du critère~$c$ pour
cette connexion :

\begin{itemize}
  \item si $c = \text{PRIX}$ → poids = prix en euros
  \item si $c = \text{TEMPS}$ → poids = durée en minutes
  \item si $c = \text{CO2}$ → poids = émissions en kg CO\textsubscript{2}e
\end{itemize}

Cette construction garantit que minimiser le poids total d'un chemin revient
exactement à minimiser le critère souhaité par l'utilisateur.

\noindent\textbf{Algorithme de résolution}

Une fois le graphe construit, on applique l'algorithme des $k$ plus courts chemins
(KPCC) en appelant \ic{AlgorithmeKPCC.kpcc(graphe,\,s,\,t,\,k)}. Cet algorithme
retourne une liste ordonnée d'au plus~$k$ objets \ic{Chemin}, chacun contenant la
liste des arêtes empruntées et le poids total du chemin. Les chemins sont sans
cycle et triés par poids croissant, ce qui correspond directement aux $k$ meilleurs
itinéraires selon le critère de l'utilisateur.

Cette modélisation est pensée pour permettre une évolution naturelle vers la
Version~2 : il suffira d'enrichir le graphe avec des nœuds de correspondance pour
prendre en charge les changements de modalité entre deux étapes d'un même voyage.

% ──────────────────────────────
\section{Prototype pour la Version~1}
% ──────────────────────────────

La classe \ic{GraphesV1} est une classe Java autonome et exécutable. Conformément
aux exigences du sujet, elle ne dépend d'aucune autre classe du projet, à
l'exception des interfaces \ic{Lieu} et \ic{Connexion} de la bibliothèque fournie,
qu'elle implémente sous forme de classes anonymes directement dans la méthode
\ic{main}.

La classe reproduit le scénario décrit en Section~1.1 : elle instancie les cinq
villes comme sommets, ajoute les connexions TRAIN avec leurs prix comme poids
d'arêtes, puis appelle \ic{AlgorithmeKPCC.kpcc} pour calculer les quatre meilleurs
itinéraires de Lille à Crépy-en-Valois. Les résultats sont affichés de façon
lisible, avec pour chaque itinéraire le rang, le coût total en euros et le détail
des étapes ville par ville.

\noindent\textbf{Lien vers le code source (commit avant la date limite) :}\\*
\href{https://gitlab.univ-lille.fr/sae2.01-2.02/2026/E5/-/blob/main/src/main/java/sae/transport/comparison/graphs/GraphesV1.java?ref_type=heads}{\ic{GraphesV1.java} sur GitLab (équipe E5)}

\end{document}
